% ==========================================
%  style/theorems.tex  —  定理环境与样式
%  基于 amsthm 宏包
% ==========================================

% --- 主题颜色 ---
\definecolor{winered}{rgb}{0.5,0,0}
\definecolor{structurecolor}{RGB}{122,122,142}
\definecolor{main}{HTML}{3D445F}
\definecolor{second}{HTML}{627581}
\definecolor{third}{HTML}{9D8798}
\definecolor{tianlan}{RGB}{102,179,255}

% --- 定理样式 ---
% 定义样式
\newtheoremstyle{defstyle}
	{3pt}{3pt}
	{\kaishu}
	{-3pt}
	{\bfseries\color{main}}
	{}
	{0.5em}
	{\indent 【\thmname{#1} \thmnumber{#2}】 \thmnote{(#3)}}

% 定理样式
\newtheoremstyle{thmstyle}
	{3pt}{3pt}
	{\kaishu}
	{-3pt}
	{\bfseries\color{second}}
	{}
	{0.5em}
	{\indent【\thmname{#1} \thmnumber{#2}】\thmnote{(#3)}}

% 问题样式
\newtheoremstyle{prostyle}
	{3pt}{3pt}
	{\kaishu}
	{-3pt}
	{\bfseries\color{third}}
	{}
	{0.5em}
	{\indent【\thmname{#1} \thmnumber{#2}】 \thmnote{(#3)}}

% --- 定理环境定义 ---
\theoremstyle{thmstyle}
\newtheorem{theorem}{定理}[chapter]

\theoremstyle{defstyle}
\newtheorem{definition}{定义}[chapter]
\newtheorem{lemma}[theorem]{引理}
\newtheorem{corollary}[theorem]{推论}

\theoremstyle{prostyle}
\newtheorem{proposition}{问题}[chapter]
\newtheorem{example}{例题}[chapter]
\newtheorem{remark}[theorem]{注}

% --- 自定义定理引用 ---
\newcommand{\mytheoremdef}[2]{%
	{\color{second}【#1】}%
	#2%
}

% --- 证明与解答环境 ---
\renewenvironment{proof}[1][证明]{%
	\par\heiti{#1.}\;\songti
}{%
	\qed\par
}

\newenvironment{solution}{%
	\par\heiti{解.}\;\fangsong
}{%
	\qed\par
}

% --- 文献引用标注 ---
\newcommand{\intro}[1]{%
	\rightline{\parbox[t]{5cm}{\footnotesize\fangsong\quad\quad #1}}%
}

% --- 自定义可计数彩色盒子 ---
\newcounter{mybox}
\newenvironment{mybox}[1]{%
	\refstepcounter{mybox}%
	\begin{tcolorbox}[
		colback=blue!5!white,
		colframe=blue!65!white,
		title={\centering Box \themybox: #1},
		breakable
	]%
}{%
	\end{tcolorbox}%
}

\newcommand{\myboxref}[1]{%
	\hyperref[#1]{\textbf{Box~\ref*{#1}}}%
}

% --- 带圈数字 ---
\newcommand{\circledtext}[2]{%
	\protect\tikz[baseline=(char.base)]{%
		\protect\node[shape=circle,draw=tianlan,fill=#1,inner sep=2pt,text=white] (char) {#2};%
	}%
}
